% Real-Time ROS 2 Camera Preprocessing Accelerator -- 3-5 minute talk
% Build:  latexmk -pdf slides.tex   (or)   pdflatex slides.tex
% Notes:  uncomment one of the \setbeameroption lines below to render
%         the speaker-notes PDF in place of (or alongside) the slides.

\documentclass[aspectratio=169,10pt]{beamer}

% ---- theme: prefer metropolis when installed, else fall back gracefully ----
\IfFileExists{beamerthememetropolis.sty}{%
  \usetheme{metropolis}%
}{%
  \usetheme{default}%
  \usecolortheme{seahorse}%
  \setbeamertemplate{navigation symbols}{}%
  \setbeamertemplate{footline}[frame number]%
}

% ---- speaker notes ----
% \setbeameroption{show notes}                          % notes after each slide
% \setbeameroption{show only notes}                     % notes-only PDF
% \setbeameroption{show notes on second screen=right}   % two-up handout

\usepackage[T1]{fontenc}
\usepackage[utf8]{inputenc}
\usepackage{amsmath, amssymb}
\usepackage{graphicx}
\usepackage{booktabs}
\usepackage{array}
\usepackage{xcolor}

\graphicspath{{../.asserts/}}

% ---- compact typography for dense slides ----
\setbeamerfont{footnote}{size=\tiny}
\newcommand{\code}[1]{\texttt{\small #1}}
\newcommand{\hi}[1]{\textbf{#1}}

\title{Real-Time ROS\,2 Camera Preprocessing Accelerator}
\subtitle{Bilinear Resize + Map-based Rectify on PYNQ-Z2 (Zynq-7020)}
\author{Jason Pan \and Yifei Feng}
\institute{HLS-based FPGA IP for the AMD ROS\,2 Perception Node}
\date{May 2026}

\begin{document}

% =====================================================================
\begin{frame}
  \titlepage
\end{frame}
\note{%
Hi -- our final project is two Vitis HLS image-processing kernels that
sit in front of the AMD ROS\,2 Perception Node: a bilinear
\code{resize\_kernel} and a map-based \code{rectify\_kernel}, both
targeting the PYNQ-Z2 Zynq-7020 board.  Both are verified
end-to-end through Python golden, HLS C-simulation, C-synthesis, and
RTL co-simulation.  I'll spend the next four minutes on the motivation,
each kernel's hardware schedule, and the verification numbers.%
}

% =====================================================================
\begin{frame}{Motivation: Why offload preprocessing to the PL?}
  \begin{columns}[T,onlytextwidth]
    \begin{column}{0.50\textwidth}
      \begin{itemize}\setlength{\itemsep}{0.35em}
        \item Robotics perception (detection / SLAM / VO) starts with
              \hi{resize + rectify}, every frame, every pixel.
        \item These ops are \hi{regular, per-pixel, frame-rate-critical}
              $\Rightarrow$ ideal hardware-acceleration target.
        \item AMD's ROS\,2 Perception Node uses \hi{exactly} this
              resize+rectify pattern -- realistic robotics use case.
        \item Goal: free the PS from per-pixel work so the CPU can focus
              on higher-level planning / inference.
      \end{itemize}
    \end{column}
    \begin{column}{0.48\textwidth}
      \centering
      \includegraphics[width=\linewidth,height=0.75\textheight,keepaspectratio]{Overview.png}\\
      {\tiny System overview (\code{.asserts/Overview.png})}
    \end{column}
  \end{columns}
\end{frame}
\note{%
The motivation is straightforward: any robot that does perception has a
camera-preprocessing stage in front of detection or SLAM, and that
stage -- resize and rectify -- is the same arithmetic on every pixel of
every frame.  That's exactly the workload that benefits from streaming
hardware: highly regular, easy to pipeline, hard for a CPU to keep up
with at high frame rates.  We picked this target because AMD's own
ROS\,2 Perception Node uses these two stages as PL-accelerated blocks,
so it's a realistic, validated use case rather than a toy benchmark.%
}

% =====================================================================
\begin{frame}{IP Architecture: PS/PL split, streaming fast path}
  \begin{columns}[T,onlytextwidth]
    \begin{column}{0.46\textwidth}
      \begin{itemize}\setlength{\itemsep}{0.3em}
        \item \hi{PS} stages frame + rectify maps in DDR; writes ctl regs.
        \item \hi{PL} modules:
          \begin{itemize}\setlength{\itemsep}{0.1em}
            \item frame ingress (DMA $\to$ AXI4-Stream)
            \item \code{resize\_kernel} (streaming)
            \item \code{rectify\_kernel} (\code{m\_axi} baseline)
            \item frame egress (AXI-Stream $\to$ DMA)
            \item control / status registers
          \end{itemize}
        \item Interfaces: \hi{AXI4-Lite} for control, \hi{AXI4-Stream}
              between compute blocks, \hi{m\_axi} for rectify's image
              and remap tables.
        \item Modular $\Rightarrow$ each stage verified \hi{independently}
              against its own Python golden.
      \end{itemize}
    \end{column}
    \begin{column}{0.52\textwidth}
      \centering
      \includegraphics[width=\linewidth,height=0.80\textheight,keepaspectratio]{Pipeline.png}\\
      {\tiny PS--PL partition (\code{.asserts/Pipeline.png})}
    \end{column}
  \end{columns}
\end{frame}
\note{%
The architecture is a standard PS/PL split.  The processor stages the
input frame and the rectify maps in DDR and writes the dimensions and
buffer pointers into the IP through AXI4-Lite.  Inside the PL we have
five modules: a frame-ingress DMA, the resize kernel, the rectify
kernel, an egress DMA, and a small control block.  Resize is fully
streaming -- AXI4-Stream in, AXI4-Stream out -- and rectify is a
memory-mapped baseline with four \code{m\_axi} masters.  Crucially the
two compute kernels are decomposed so we can verify each one bit-exact
against its own Python golden before we ever wire them together.%
}

% =====================================================================
\begin{frame}{Resize Kernel: bilinear, 2-row ring buffer, II=1}
  \begin{columns}[T,onlytextwidth]
    \begin{column}{0.52\textwidth}
      \textbf{Math (\code{float32}, matches Python golden bit-exact):}
      \[
        u = s_x x,\quad v = s_y y,\qquad
        s_x = \tfrac{W_\text{in}}{W_\text{out}},\;
        s_y = \tfrac{H_\text{in}}{H_\text{out}}
      \]
      \[
        I_\text{out}(x,y)=\sum_{i,j\in\{0,1\}} w_{ij}\,I_\text{in}(u_i,v_j)
      \]
      \textbf{Schedule -- key insight:} $v(y)$ is monotonic in $y$
      $\Rightarrow$ only \hi{two consecutive input rows} are ever live.
      \begin{itemize}\setlength{\itemsep}{0.2em}
        \item \code{pixel\_t line\_buf[2][MAX\_W]} (no full frame in BRAM)
        \item \code{ARRAY\_PARTITION dim=1 complete}: both taps in 1 cycle
        \item \code{BIND\_STORAGE ram\_t2p impl=bram}: read + write parallel
        \item \code{PIPELINE II=1} on \code{LOAD\_ROW} \& \code{EMIT\_PIXEL}
      \end{itemize}
    \end{column}
    \begin{column}{0.45\textwidth}
      \footnotesize
      \begin{verbatim}
for r in 0..in_h-1:
  LOAD_ROW: stream row r       // II=1
            -> line_buf[r&1]
  EMIT: while v1(next_y) <= r:
    EMIT_PIXEL:                // II=1
      for xo in 0..out_w-1:
        4 taps from line_buf
        out_stream <- bilinear
\end{verbatim}
      \vspace{0.4em}
      \hi{Steady state: 1 pixel / cycle.}\\
      Single read of each input row, single write of each output pixel.
    \end{column}
  \end{columns}
\end{frame}
\note{%
The resize kernel is where the cleanest hardware story is.  The math is
plain float32 bilinear interpolation; the trick is the schedule.  Since
$v$ is monotonic in $y$, the only input rows you ever need are
\code{v0} and \code{v0+1} -- two rows total, not the whole frame.  We
hold them in a 2-row line buffer in BRAM, fully partition the
2-row dimension so both 2$\times$2 taps come out in the same cycle, and
bind that buffer to a true-dual-port BRAM so the load side and the
emit side can run concurrently.  With \code{PIPELINE II=1} on both
inner loops, the kernel sustains \hi{one pixel per cycle} in steady
state, with each input row touched once and each output pixel written
once -- no backpressure, no slow path.%
}

% =====================================================================
\begin{frame}{Rectify Kernel: \code{m\_axi} baseline + optimisation roadmap}
  \begin{columns}[T,onlytextwidth]
    \begin{column}{0.55\textwidth}
      \textbf{Math (reuses resize's bilinear sampler):}
      \[
        u = \mathrm{map}_x[y,x],\quad v = \mathrm{map}_y[y,x]
      \]
      \[
        I_\text{rect}(x,y)=\sum_{i,j}w_{ij}\,I_\text{in}(u_i,v_j)
      \]
      \textbf{Baseline interface:} four parallel \code{m\_axi} masters
      \begin{itemize}\setlength{\itemsep}{0.15em}
        \item \code{gmem\_img} \,\, $\to$ image (4 taps / px)
        \item \code{gmem\_mapx}, \code{gmem\_mapy} $\to$ float remap tables
        \item \code{gmem\_out} \, $\to$ output image
      \end{itemize}
      Memory-bound: HLS settles on \hi{II = 4 cyc/px}.
    \end{column}
    \begin{column}{0.43\textwidth}
      \textbf{Measured performance (1920$\times$1080, PYNQ-Z2 @100 MHz):}\\[0.3em]
      \begin{tabular}{ll}
        \toprule
        cycles       & 8{,}294{,}471 \\
        wall @100MHz & $\approx$ \hi{82.9 ms} ($\approx$ 12 fps) \\
        \bottomrule
      \end{tabular}
      \vspace{0.6em}
      \textbf{Roadmap to II=1 (deferred, in \code{plan.md}):}
      \begin{itemize}\setlength{\itemsep}{0.1em}
        \item line buffer of depth $D_y{+}2$ from offline map bounds
        \item schedulable output queue, in-window hit test
        \item DDR slow-path fallback for rare misses
      \end{itemize}
      Honest scope: baseline shipped \& verified; optimised path designed
      but not implemented.
    \end{column}
  \end{columns}
\end{frame}
\note{%
Rectify reuses the same float32 bilinear sampler, but coordinates come
from two external remap tables instead of a fixed scale factor.  We
shipped this as an intentionally simple baseline: four parallel
\code{m\_axi} masters -- image, \code{map\_x}, \code{map\_y}, output --
no on-chip line buffer.  HLS converges on II=4 cycles per pixel,
bottlenecked by memory dependencies.  On the PYNQ-Z2 at 100\,MHz the
csynth report measures 8.29\,M cycles, about 82.9\,ms for a full HD
frame -- roughly 12\,fps, well under 30.  The
path to II=1 is in our plan: line-buffer depth $D_y{+}2$ derived from
an offline analysis of the map, an output queue that defers pixels
until their source rows are resident, and a DDR slow-path for misses.
We documented it but did not implement it -- scope was honest.%
}

% =====================================================================
\begin{frame}{Verification: golden-first, four layers, zero tolerance}
  \begin{columns}[T,onlytextwidth]
    \begin{column}{0.50\textwidth}
      \begin{enumerate}\setlength{\itemsep}{0.3em}
        \item \hi{Python golden} (\code{sw/resize.py}, \code{sw/rectify.py})
              -- \code{float32}, unit-tested vs hand-computed,
              \code{numpy} vectorized, OpenCV cross-check optional.
        \item \hi{HLS C-simulation} -- C++ port of golden with the
              \emph{same} float32 op-order $\Rightarrow$ \code{memcmp}
              with \hi{no tolerance window}.
        \item \hi{C-synthesis} -- resource \& schedule reports
              (\code{csynth.rpt}).
        \item \hi{C/RTL co-simulation} -- runs the same testbench
              through generated Verilog.
      \end{enumerate}
    \end{column}
    \begin{column}{0.48\textwidth}
      \textbf{Test corpus:}
      \begin{itemize}\setlength{\itemsep}{0.15em}
        \item Resize: 5 synthetic shapes + 5 NYU real images
              (downscale, upscale, identity, non-integer, wide).
        \item Rectify: 5 deterministic cases
              (identity, fractional shift, crop, border clamp, scale-map).
        \item All comparisons are \hi{byte-exact} against
              \code{\_gold.bin} from the Python vectorised model.
      \end{itemize}
      \vspace{0.3em}
      \footnotesize{Real-data fixtures live under
      \code{hls/testdata/}, regenerated by
      \code{sw/prepare\_real\_data.py}.}
    \end{column}
  \end{columns}
\end{frame}
\note{%
Verification is the core methodology story.  We write a Python golden
first, unit-test it against hand-computed 2$\times$2 examples and
against OpenCV when it's available.  Then the HLS C++ kernel is a
line-by-line port of that golden, keeping the float32 op-order
identical so the testbench can \code{memcmp} outputs with zero
tolerance.  After C-simulation we run C-synthesis for resource and
schedule reports, then C-RTL co-simulation to confirm the generated
Verilog matches.  The corpus is five synthetic resize cases plus five
NYU indoor images, and five deterministic rectify cases that exercise
identity, fractional shift, cropping, border clamping, and non-integer
scaling -- all byte-exact against the golden binaries.%
}

% =====================================================================
\begin{frame}{Results on \code{xc7z020} @ 10\,ns: functional green, resize needs clock relax}
  \textbf{Verification:}\\[0.3em]
  \centering
  \begin{tabular}{lcccc}
    \toprule
    kernel  & py-tests & C-sim & csynth slack & co-sim \\
    \midrule
    resize  & 5/5 & 5/5 & $-0.70$\,ns (closes $\approx$93\,MHz) & \hi{10/10} (5 synth + 5 NYU)\\
    rectify & 4/4 & 5/5 & $0.00$\,ns (clean)                    & 5/5 \\
    \bottomrule
  \end{tabular}\\[0.7em]

  \textbf{Resource utilisation (from \code{csynth.rpt}):}\\[0.3em]
  \begin{tabular}{lcccccc}
    \toprule
    kernel  & BRAM & DSP & FF & LUT & \% of part (LUT/DSP) & throughput \\
    \midrule
    resize  & 2  & 37 & 4{,}852 & 8{,}699 & \hi{16\,\% / 16\,\%} & 1 px / cycle \\
    rectify & 12 & 13 & 7{,}508 & 9{,}623 & \hi{18\,\% / 5\,\%}  & 4 cyc / px   \\
    \bottomrule
  \end{tabular}\\[0.3em]
  {\tiny For reference, on the larger \code{xczu7ev} @ 5\,ns each kernel sat
   at 1--3\,\% on every metric -- the jump reflects \code{xc7z020}'s
   smaller LUT/DSP budget, not a design regression.}\\[0.5em]

  \raggedright
  \textbf{Takeaway:} Functional verification carried cleanly to PYNQ-Z2;
  rectify meets 100\,MHz, resize needs a slightly relaxed clock
  ($\approx$90\,MHz) on \code{xc7z020} -- a normal port adjustment.  The
  streaming line-buffer schedule keeps resize compact; rectify is
  memory-bound until the buffered path lands.
\end{frame}
\note{%
Final numbers from the re-synthesis on \code{xc7z020} at 10\,ns:
every functional verification layer still passes for both kernels --
the four-layer flow doesn't depend on the part, so porting from
\code{xczu7ev} to PYNQ-Z2 broke nothing we'd already proved.
Resource-wise, the design is no longer 1--3\% of the part; it's
roughly 16\% LUT for resize and 18\% LUT for rectify, because
\code{xc7z020} is a much smaller device.  The one real issue is that
the resize \code{EMIT\_PIXEL} pipeline reports $-0.70$\,ns of slack at
the 10\,ns target -- it closes at about 93\,MHz, so Vivado P\&R will
need a slightly relaxed clock for sign-off.  Rectify meets 100\,MHz
cleanly and lands its full HD frame in about 83\,ms.  The qualitative
takeaway holds: streaming schedule keeps resize compact, rectify
needs the buffered variant for higher fps.  Happy to take questions.%
}

\end{document}
